% One Britannica scan as extracted by Nemotron Parse 3.5. Filled by src/nemotron_challenge/stages/latex.py.
% The TITLE, BODY and SOURCE markers (each wrapped in two percent signs) are replaced. Compile with XeTeX (tectonic).
\documentclass[10pt]{article}
\usepackage[a4paper,margin=18mm]{geometry}
\usepackage{fontspec}
\usepackage{multicol}
\usepackage{graphicx}
\usepackage[hidelinks]{hyperref}
\setlength{\parindent}{0pt}
\setlength{\parskip}{4pt}
\setlength{\columnsep}{8mm}
\begin{document}
\section*{Historical source: PHILLOOGY}
\begin{multicols}{2}
\subsection*{\textit{continued:} PHILLOOGY}
Cushim was the same. We may therefore rest assured, that whatever changes the Ethiopian dialect may have un-dergone in the course of three thousand years, it was originally either Chahidae, or at least a branch of that language. Sadger informs us, that the Ethiopian call themselves Chadleans; and this, says he, not without reason, because of the largest secret and prende books which are extant amongst them, the next elegant and most beautiful and written in a style resembling that of the Chaldeau or Asyrian: Masani Victories and then that reduced the Ethiopians to go to the rules of grammar, tells an in his \emph{Procommon}, that the Ethiopia can tell their tongue Chadlars, that it springs from the Babylonian, and is very like the Hebrew, Syria, and Arabic; and at the same time, be concludes, that this language may be easily learned by those who no masters of the Hebrew. The learned there, and Baby Walton in his \emph{Prolegemaa}, are clearly of the same opinion. The writer letters of the Ethiopians, according to Dio-dorus Sicula, were the same with the sacred characters of the Egyptians." From this account, if the Sicilian has be reached, the sacred letters of these people, concerning which so many vague conjectures have been formed, were actually Chahidio. To cover on this investigation a little matter, we may observe, that Six Villan Jones seems to have paved, by very plausible arguments, that the Sanscrit character were deduced from the Chahida. This circumcments affords a presumption that the Ethiopian Cushin were likewise concerned with the Egyptians; who, as is remarked in the section concerning the \emph{Snowsic}, probably introduced the religion of the Brahman into Hisolution. This is advanced as a conjecture only; and yet when we consider the affinity between the Egyptian and the Grenzo religions, are strongly inclined to believe that this simism may one day for its verbularable facts. The original Ethiopians were a highly civilized people; and their laws, their institutions, and especially their religion, and celebrated far and wide. Hence talks to physics as the piety of the Ethiopians, and sends his gods every now and then to read twelve days with that devout people. The Sciland produces a number of very speakers arguments to prove that those two nations equally sprung from the same stock. He mentions a similarity of features, of manners, of anatomy of law, of letters, of the fashion of statuses, and of religion, as evidences of the relation between these two neighboring nations. There was, since every body knows a common, in regard to sacrifice it, even the two countries. The Egyptians sent annually a deputation of their priests, limited with the probable nature of their gods, to visit the fines of the devon Ethiopians. Upon this occasion, a solemal religious banque was prepared, which lasted twelve days, and of which the priest of both nations were parties. It was a kind of sacramental institution, we imagine, by which both parties publicly awarded their agreement in the creation of their religion respectively. These observations plainly show, that the most ancient Ethiopians were people highly utilized in need to much, but the Egyptians were all one time connected to be their scholars. The tone of their language was certainly the same with that of the Chaldeau or Arabica Cunha, whom whom they are descended. We know not whether there are any books in the ancient Ethiopians now act; so that it is not easy to procide instance of its coincidence with the Chaldaic. But Diogenes Lacertas’ Roman us, that Tarnayihus, in his catalogue of the books composed by Democritus, mentions one, 1 Lib. 10. p. 10. j. Sept. We also the same observation confirmed by Helodorous (\emph{Ethiopia, lib} x. p. 476.) " The royal letters of the Ethiopians," says io. \emph{regi ruse le Mopoe oppaar oppaarna,} concerning the sacred last. Chahidic ties in the island of Morce 1 and another concerning the Language, sacred therein. Had these books survived the ravages of time, they would, in this age of research and consisary, have determined not only the point under the consideration, but likewise the affinity of sacred rites amongst the Chaldeana, Ethiopians, and Egyptianman. We knew now down that the Ethiopians were a colony Modern of Cushines; that the Cushies were originally sovereigns of Ethiopie Shame and Charles, and consequently spoke other Chaldaic, \emph{burgues} - or a dialect of the tongue that their colonies must have used the same languages; and that the ancient Ethiopians were a people highly published, and celebrated in the next early ages on account of their virtue and piety. It has ever appeared that the citizens letters of that people were the sacred characters of the Companies. These letters, we incing, were the Cupile (; for some account ofwhich see the section on the Arabic). When the way were descendants, and the modern letters substituted in their room, cannot be determined; nor is it, we apprehend, a matter of much importance. We shall therefore say that part of the subject, and refer our curious and impulsive readers to the favored Jod Ladol's excellent grammen and dictionary of the Abyssian or ees tongue; where they will find every thing worth knowing or that subject. We shall endeavour to grafity on readers with a very brief account of the modern Ethnic of Avisosion tonges, few whatso we are principally indebted to that learned, indetatigable, and adventurous travelice. Bruce , who, by his observations on that country, which he made in person, often at the board of hisft, has discovered, as it were, a new world both to Europe and Asia. The rare ancient language of Ethiopians, which we shall now call by its modern name of Abyssian, was according to that gentleman, the Geest, which was spoken by the ancient Chinese shopdesks. This was should think a approaches toerst to the old Chaldaic. Upon a revolution in Haheals, the court resided many years in the province of Anhara, where the people spoke a different language, or at least a very different dialect of the same language. During this interval, the Geese, or language of the subjects, was dropped, and redefined only in writing, as a dead language; the sacred Scriptures being in that tongue and saved it from going line across. This dialect is, escending had and unamomnion. It is full of the letters, I and Y, on which an accent is just that nearly resembles something. Considering the small town of sea which divides this countr try from Arabia, we need not wonder that it has great affinity with the Arabic. It is not difficult to be captured by those who understand any other or certain languages; and as the roots of many Hebrew words are only to be found here, lectures to be absolutely necessary to all those who wish to obtain a critical skill in that language. The Ethiopie alphabet consists of twenty-six letters, each Ethiopie of which, by a stripped or point emerged, writes its second alphabet. In such a manner as that those twenty-six forms as it were sixty-two distinct letters. At first they had only seventy-five of these original letters, the Ladr P being wanting; so that they were obliged to substitute another letter in its place. \_Ethikale, the example, they call Pauline, Aruba, or Guina, and Peteror they pronounced Keten.\_ At six they substituted T, and added this to the end of their alphabet; giving up the force of B, though it was really in expectation of an old character, rather than in inventof a war one. Besides these, there are twenty others of the nature of diph were the sacred characters of the Egyptians." Casslodorus likewise sesurs us, "that the letters inserted upon the Egyptian cheliska were the Palestry." Intentery discovers have refused these and other similar statements. 2 Lib. in de 401., Comisiones. Where the capital of Ethiopia was situated. A very learned German, who published a grammar and dictionary of the Geez in folio-
\end{multicols}
\section*{Source record}
\begin{itemize}
\item Page ID: \texttt{britannica:309f765b1ecc83e0a7f7762e44b4a927a7555443:train:403}
\item Scan: \texttt{page-0403-44f912e1c323.jpg}, 600 x 823 px
\item Image SHA-256: \texttt{\small 44f912e1c3237b35750b6e188e94f3a20f2af3f1dba0a4bdb630dff640ddbf3c}
\item Extraction: \texttt{nvidia/nemotron-3.5-parse}, raw reply \texttt{artifacts/extract/raw/0403.json}
\item Text is Parse OCR output, normalized (long s as s, line-end hyphens joined). It is not corrected by hand.
\end{itemize}
\end{document}
